\documentclass{article}

% Use the Cactus ThornGuide style file
% (Automatically used from Cactus distribution, if you have a 
%  thorn without the Cactus Flesh download this from the Cactus
%  homepage at www.cactuscode.org)
\usepackage{../../../../doc/latex/cactus}

\begin{document}

\title{CottonmouthBSSN}
\author{Lucas Sanches \\ Max Morris \\ Steven R. Brandt}
\date{$ $Date$ $}

\maketitle

% Do not delete next line
% START CACTUS THORNGUIDE

% Add all definitions used in this documentation here

\begin{abstract}
    Provides the BSSN equations as generated by the EinsteinEngine.
\end{abstract}

\section{Introduction}

%Einstein Engine
The Einstein Engine is a DSL and toolkit for creating Cactus thorns.
It can generate complete CarpetX thorns from simple recipes written in Python.

\subsection{How to set up EinsteinEngine}

\begin{itemize}
   \item Create a venv and activate it. 
   \begin{verbatim}
   python -m venv venv && . ./venv/bin/activate
   \end{verbatim}
\item Install the dependencies.
   \begin{verbatim}
   python -m pip install -r requirements.txt
   \end{verbatim}

\item Install EinsteinEngine. 
   \begin{verbatim}
   python -m pip install .
   \end{verbatim}
\end{itemize}
   
\subsection{How to generate a thorn}

\begin{itemize}
\item Be sure the venv is activated.
   \begin{verbatim}
   . ./venv/bin/activate
   \end{verbatim}

\item Typecheck the recipe with MyPy.
   \begin{verbatim}
   mypy recipes/Cottonmouth/bssnok_unsplit_ricci_opt.py
   \end{verbatim}

\item Run the recipe.
   \begin{verbatim}
   python recipes/Cottonmouth/bssnok_unsplit_ricci_opt.py
   \end{verbatim}
\end{itemize}

\section{Parameters}

TBD

% Do not delete next line
% END CACTUS THORNGUIDE

\end{document}
